% §4 The Latency Floor Theorem
We now prove the central result: the silence gap of any cascade pipeline has an
irreducible lower bound.

\subsection{The Latency Floor Theorem}

\begin{theorem}[Latency Floor]
\label{thm:latency_floor}
Let $G = (S, E)$ be an \emph{EOU-gated cascade} (Definition~\ref{def:egc}) with stages
$S = \{s_1, \ldots, s_n\}$ forming a linear chain.
Let each stage $s_i$ have latency decomposed as $\lambda_i = \lambda_i^r + \lambda_i^{irr}$
per Definition~\ref{def:decomposition}.
The following premises are assumed:
\begin{enumerate}
  \item \emph{Sequential execution}: stages execute strictly sequentially; no stage begins
        processing before its predecessor has produced its complete output.
        (This premise holds by definition for EGC; it does \emph{not} hold for the streaming
        cascade of Definition~\ref{def:sc}, which is analysed separately in
        Theorem~\ref{thm:streaming_floor}.)
  \item \emph{Input independence}: each stage latency $\lambda_i$ is treated as independent of the specific input $x_i$ (a simplifying assumption acknowledged in Section~\ref{sec:discussion}).
  \item \emph{Implementation}: an ``implementation'' of $G$ is any choice of model weights, inference hardware, and serving configuration that preserves the pipeline topology and stage boundary definitions; the set of valid implementations for each stage is $\mathrm{valid}(s_i)$ as defined in Definition~\ref{def:decomposition}.
\end{enumerate}
Then the silence gap satisfies:
\begin{equation}
  \Lambda \geq \Lambda^* := \sum_{i=1}^{n} \lambda_i^{irr}
\end{equation}
for all implementations of $G$, regardless of model size, hardware, or inference optimisation.
\end{theorem}

\begin{proof}
Each stage $s_j$ in the linear chain cannot begin execution until $s_{j-1}$ produces its output,
by Definition~\ref{def:cascade_dag} (Cascade Pipeline DAG).
Therefore the silence gap is at least the sum of the first-output latencies of all stages:
$\Lambda \geq \sum_i \lambda_i$.
Since $\lambda_i = \lambda_i^r + \lambda_i^{irr}$ and $\lambda_i^r \geq 0$,
we have $\Lambda \geq \sum_i \lambda_i^{irr} = \Lambda^*$.
The reducible terms $\lambda_i^r$ can be made arbitrarily small by using faster models
(in the limit of infinite compute), but $\lambda_i^{irr}$ is by definition the lower bound
achievable for any implementation of $s_i$ given its input--output contract.
Therefore $\Lambda^*$ is a lower bound on $\Lambda$ for any cascade pipeline with topology $G$.
\end{proof}

\begin{remark}
Theorem~\ref{thm:latency_floor} applies to any linear or branching DAG topology.
For a branching pipeline where parallel stages can execute concurrently,
$\Lambda^*$ is the maximum irreducible latency along any critical path through the DAG,
not the sum over all stages.
The standard cascade pipeline is a single critical path,
so the sum form applies directly.
\end{remark}

\subsection{Native Voice as the Zero-Floor Limit}

A native voice model collapses the three-stage cascade into a single stage $s_{native}$
that maps audio input directly to audio output, eliminating inter-stage buffering latency.
Two architecturally distinct classes exist, with different prefill-hiding properties.

\paragraph{Buffer-and-commit (GPT-4o Realtime class).}
Systems such as GPT-4o Realtime \citep{openai2024realtimeapi} and Gemini Live
\citep{google2024geminivoice} stream audio into a server-side buffer during the user's speech.
VAD monitors the stream and, upon detecting end-of-utterance, commits the buffer and triggers
model inference.
The full audio-token prefill therefore begins \emph{after} EOU, analogously to EGC cascade:
\begin{equation}
  \Lambda^*_{RT,commit} = d_{eou}^{vad} + \alpha_{audio} \cdot N_{speech} + \lambda_{decode,1}
\end{equation}
The advantage over cascade is the elimination of the separate ASR decode and TTS stages,
not prefill hiding during speech.
Observed silence gaps of 320--500~ms are consistent with this structure.

\paragraph{Streaming causal (Moshi class).}
Systems such as Moshi \citep{defossez2024moshi} use a fully causal audio codec
(frame stride $\Delta_t = 80$~ms) and process each codec frame through the Temporal Transformer
as it arrives.
The KV cache grows continuously during the user's utterance:
by the time speech ends, the model has already prefilled all speech frames.
The post-EOU silence gap is therefore:
\begin{equation}
  \Lambda^*_{RT,stream} = \Delta_t + \lambda_{forward}^{(1)}
\end{equation}
where $\lambda_{forward}^{(1)}$ is the time for one Transformer step over the current KV state.
This gives a theoretical floor of approximately 160~ms, independent of utterance length.
The prefill advantage relative to EGC cascade is:
\begin{equation}
  \Lambda^*_{EGC} - \Lambda^*_{RT,stream}
    = \underbrace{\lambda_{asr}^{irr} + \lambda_{tts}^{irr}}_{\text{stage elimination}}
    + \underbrace{\alpha_{audio} \cdot N_{speech} - \lambda_{forward}^{(1)} - \Delta_t}_{\text{prefill hidden during speech}}
\end{equation}
The prefill-hiding term grows with utterance duration $T_{speech}$,
so streaming causal models gain a larger advantage for longer utterances.

The formal inequality chain is:
\begin{equation}
  \Lambda^*_{RT,stream} \leq \Lambda^*_{RT,commit} \leq \Lambda^*_{EGC}
\end{equation}
where the first gap reflects prefill hiding during speech and the second reflects
ASR and TTS stage elimination.
The paper's impossibility result (Corollary~\ref{cor:impossibility}) applies only to EGC cascade,
not to streaming cascade or native voice architectures.

\subsection{The Cascade Competitiveness Criterion}

\begin{definition}[Cascade Competitiveness Criterion]
\label{def:ccc}
A cascade pipeline is \emph{perceptually competitive} with a native voice model if and only if:
\begin{equation}
  \Lambda \leq \tau_{human}
\end{equation}
where $\tau_{human}$ is the human turn-taking threshold (approximately 200~ms,
as established by \citet{levinson2015timing} and \citet{stivers2009universals}).
\end{definition}

\begin{corollary}[Impossibility of Cascade Competitiveness]
\label{cor:impossibility}
If $\Lambda^* > \tau_{human}$, then no \emph{EOU-gated cascade} pipeline with topology $G$
can satisfy the Cascade Competitiveness Criterion, regardless of model or hardware improvements.
\end{corollary}

\begin{proof}
By Theorem~\ref{thm:latency_floor}, $\Lambda \geq \Lambda^*$.
If $\Lambda^* > \tau_{human}$, then $\Lambda > \tau_{human}$ for all implementations,
so the CCC $\Lambda \leq \tau_{human}$ cannot be satisfied.
\end{proof}

\begin{remark}
Corollary~\ref{cor:impossibility} is scoped to EOU-gated cascade.
Streaming cascade (Definition~\ref{def:sc}) violates premise~(1) of Theorem~\ref{thm:latency_floor}
by design, and therefore has a strictly lower latency floor, derived in Theorem~\ref{thm:streaming_floor} below.
\end{remark}

\begin{theorem}[Streaming Cascade Floor]
\label{thm:streaming_floor}
Let $G$ be a streaming cascade (Definition~\ref{def:sc}).
Let $\delta_N$ be the commit latency of the final word from EOU,
$\alpha$ the per-token LM prefill cost,
$|C_{residual}|$ the number of tokens not yet prefilled at EOU,
$\lambda_{lm}^{decode,irr}$ the time to generate the first output token from a
completed KV cache, and $\lambda_{tts}^{irr}$ the TTS irreducible latency.
Then:
\begin{equation}
  \Lambda_{SC} \geq \Lambda^*_{SC}
    := \delta_N + \alpha \cdot |C_{residual}| + \lambda_{lm}^{decode,irr} + \lambda_{tts}^{irr}
\end{equation}
The gain over the EGC floor is the prefill work amortised over the user's speech duration:
\begin{equation}
  \Lambda^*_{EGC} - \Lambda^*_{SC}
    = \alpha \cdot (|C_{user}| - |C_{residual}|)
    = \alpha \cdot |C_{pre}|
\end{equation}
where $|C_{pre}|$ is the number of tokens stably committed and prefilled during speech.
\end{theorem}

\begin{proof}
At $t = 0$ (EOU): (1)~$\delta_N$ has elapsed for the ASR to commit the final word; (2)~the
LM decode phase cannot begin before EOU by Definition~\ref{def:sc}(c); (3)~the residual
$\alpha \cdot |C_{residual}|$ tokens must be prefilled before the first output token is generated;
(4)~TTS cannot begin before the first LM output token.
Each term is irreducible under its respective stage contract.
The gain $\alpha \cdot |C_{pre}|$ is bounded above by $\lambda_{lm}^{irr}$ (the full EGC
prefill cost) and approaches zero as $T_{speech} \to 0$ (short utterances leave no time for overlap).
\end{proof}

\begin{corollary}[Streaming CCC Satisfiability]
\label{cor:streaming_ccc}
A streaming cascade can satisfy the Cascade Competitiveness Criterion if and only if:
\begin{equation}
  \delta_N + \lambda_{lm}^{decode,irr} + \lambda_{tts}^{irr} \leq \tau_{human}
\end{equation}
Substituting representative values ($\lambda_{lm}^{decode,irr} \approx 20$~ms,
$\lambda_{tts}^{irr} \approx 80$~ms, $\tau_{human} = 200$~ms), this requires
$\delta_N \leq 100$~ms.
Streaming RNN-T systems \citep{sainath2020streaming} achieve $\delta_N \approx 150$--320~ms
at production accuracy, placing the CCC at the boundary of feasibility.
Fully causal CTC architectures with minimal lookahead can achieve $\delta_N \approx 100$~ms
at a 5--15\% relative word error rate cost \citep{yao2021wenet}.
\end{corollary}

Figure~\ref{fig:perceptual_threshold} illustrates the relationship between
the latency floor $\Lambda^*$, the total silence gap $\Lambda$, and the human threshold.

\begin{figure}[t]
  \centering
  \resizebox{\linewidth}{!}{%
    \begin{tikzpicture}[x=0.01354cm, y=1cm]

% CCC satisfied zone (green background)
\fill[green!8] (0,0) rectangle (200,3);

% Axis
\draw[-{Stealth}, thick] (0,0) -- (960,0)
  node[right, font=\small] {Silence gap $\Lambda$ (ms)};
\draw[-{Stealth}, thick] (0,-0.3) -- (0,3.2);

% Threshold markers
\draw[dashed, thick, red!70!black] (200,0) -- (200,3)
  node[above, font=\footnotesize, align=center] {$\tau_{human}$\\200 ms};
\draw[dashed, thick, orange!80!black] (800,0) -- (800,3)
  node[above, font=\footnotesize, align=center] {Discomfort\\onset\\800 ms};

% Native voice band (320--500 ms; lies entirely outside the CCC zone)
\fill[green!20] (320,2.0) rectangle (500,2.7);
\draw[green!60!black, thick] (320,2.0) rectangle (500,2.7);
\node[font=\tiny, align=center, text width=2.4cm] at (410,2.35)
  {Native voice\\(GPT-4o RT: 320--500 ms)};

% Cascade optimised band
\fill[blue!15] (700,1.2) rectangle (900,1.9);
\draw[blue!60!black, thick] (700,1.2) rectangle (900,1.9);
\node[font=\scriptsize, align=center, text width=2.6cm] at (800,1.55)
  {Cascade (typical)\\700--900 ms};

% Cascade floor band
\fill[red!15] (300,0.3) rectangle (500,1.0);
\draw[red!60!black, thick] (300,0.3) rectangle (500,1.0);
\node[font=\tiny, align=center, text width=2.2cm] at (400,0.65)
  {$\Lambda^*$ floor\\(best-case cascade)};

% CCC zone label (rotated)
\node[font=\footnotesize, color=green!50!black, rotate=90] at (100,1.5)
  {CCC satisfied};

% x-axis ticks
\foreach \x in {0,200,400,600,800} {
  \draw (\x, 0.05) -- (\x, -0.05)
    node[below, font=\scriptsize] {\x};
}

\end{tikzpicture}
%
  }
  \caption{Perceptual threshold diagram.
           The green zone ($\Lambda \leq \tau_{human} = 200$~ms) is the region where
           the Cascade Competitiveness Criterion is satisfied.
           Buffer-and-commit native voice models such as GPT-4o Realtime (green band,
           320--500~ms) fall outside the strict 200~ms CCC zone, yet are substantially
           below typical cascade latency.
           A typical cascade pipeline (blue band) operates well beyond it.
           The best-case cascade floor $\Lambda^*$ (red band) remains above
           $\tau_{human}$, establishing the impossibility result of
           Corollary~\ref{cor:impossibility}.}
  \label{fig:perceptual_threshold}
\end{figure}

\subsection{Practical Interpretation of \texorpdfstring{$\tau_{human}$}{tau\_human}}

The threshold $\tau_{human} = 200$~ms is derived from the modal inter-turn interval observed in human-human conversation \citep{levinson2015timing,stivers2009universals}.
We note that human tolerance for machine voices is empirically somewhat higher than for human interlocutors: users of voice assistants typically accept gaps up to 300--500~ms without explicit dissatisfaction \citep{levinson2015timing}.
The CCC should therefore be interpreted as a \emph{target} derived from the ideal human-human norm: cascade architectures that approach this threshold are perceptually competitive, and those that substantially exceed it are not.
Native voice systems such as GPT-4o Realtime, which report silence gaps of 320--500~ms \citep{openai2024realtimeapi}, themselves fall outside the strict $\tau_{human} = 200$~ms zone; they nonetheless set the practical bar for perceptual competitiveness as the current state of the art for low-latency voice interaction.
The impossibility result (Corollary~\ref{cor:impossibility}) applies in the strict sense: cascade pipelines whose irreducible floor $\Lambda^* > \tau_{human}$ cannot approach the ideal human-human norm without architectural change, and are structurally at a disadvantage relative to native voice models regardless of model or hardware improvements.

\subsection{Implications}

Corollary~\ref{cor:impossibility} has a direct engineering implication:
for the representative latency values in Table~\ref{tab:latency_example},
where $\Lambda^* = 380$~ms $> \tau_{human} = 200$~ms,
no model improvement or infrastructure upgrade can make the cascade pipeline
perceptually competitive in the strict sense.
The pipeline topology must change---through streaming synthesis, barge-in, or hybrid
routing---for the CCC to become satisfiable.
We formalise these architectural modifications in Section~\ref{sec:design_levers}.
